\documentclass[a4paper,11pt]{article}
% Single-file note (conventions as in notes/chromaticity-v1/preamble.tex)
\usepackage[T2A]{fontenc}
\usepackage[utf8]{inputenc}
\usepackage[russian]{babel}
\usepackage{amsmath}
\usepackage{amssymb}
\usepackage[hidelinks]{hyperref}
\numberwithin{equation}{section}
\setlength{\emergencystretch}{2em}
\newcommand{\rf}{\mathrm{rf}}
\newcommand{\cE}{\mathcal{E}}

\title{Гамильтониан ВЧ-резонатора в канонических переменных $(x,p_x,y,p_y,l,\delta)$}
\author{}
\date{}

\begin{document}
\maketitle

\begin{abstract}
Записан $s$-гамильтониан ускоряющего ВЧ-резонатора в переменных, где $\delta=\Delta p/p_0$~--- отклонение импульса,
а $l=-v\,\Delta t$~--- продольная координата, канонически сопряжённая $\delta$. Гамильтониан выводится по шагам
непосредственно в этих переменных. Рассмотрены поле моды TM$_{010}$, усреднение по длине (фактор пролётного времени),
тонкий предел и его точность (формула Бейкера--Кэмпбелла--Хаусдорфа), симплектичность, различие накопителя и
ускорения (автономность, смена опорного импульса, адиабатическое затухание), частота синхротронных колебаний $Q_s$,
продольные уравнения движения и точное однооборотное отображение в переменных $(l,\delta)$ и $(\varphi,\Delta E)$
с разложением коэффициента проскальзывания до второго порядка ($\eta_1,\eta_2$). Связь с переменными
$(z,\delta_E)$ из лекций~\cite{wolski4,wolski7}~--- в приложении.
\end{abstract}

\section{Переменные и гамильтониан}
Опорная частица: импульс $p_0$, скорость $v_0=\beta_0c$, полная энергия $E_0=\gamma_0mc^2$, $\beta_0\gamma_0=p_0/(mc)$;
заряд частицы $Ze$. Независимая переменная~--- длина опорной орбиты $s$. Опорная частица проходит точку $s$
в момент $t_0(s)=\int ds/v_0$ (при постоянном $v_0$ это $s/v_0$); \emph{запаздывание} частицы
\[
\Delta t = t - t_0(s)
\]
($\Delta t<0$~--- частица пришла раньше опорной). Частица имеет импульс $p$, скорость $v=\beta c$, полную энергию
$E=\gamma mc^2=\sqrt{p^2c^2+m^2c^4}$. Определим две продольные пары:
\begin{align}
&(q,P):\qquad q=-c\,\Delta t,\qquad P=\frac{E}{p_0c}, \label{eq:qP}\\
&(l,\delta):\qquad l=-v\,\Delta t,\qquad \delta=\frac{p-p_0}{p_0}. \label{eq:ld}
\end{align}
Здесь $q$~--- опережение частицы, выраженное в длине (путь, проходимый светом за $-\Delta t$), $P$~--- полная энергия
в единицах $p_0c$; $l$~--- опережение в длине, но с собственной скоростью частицы (расстояние, на которое частица
впереди опорной, если в момент прихода опорной частицы она уже прошла это расстояние), $\delta$~--- относительное
отклонение импульса. Поперечные импульсы $p_{x,y}=P_{x,y}/p_0$ (канонические, нормированные на $p_0$).
Удобно ввести безразмерную полную энергию как функцию $\delta$:
\begin{equation}
\cE(\delta)=\sqrt{(1+\delta)^2+\Big(\frac{mc}{p_0}\Big)^2}=\frac{E}{p_0c},\qquad
\Big(\frac{mc}{p_0}\Big)^2=\frac{1}{\beta_0^2\gamma_0^2},\qquad \cE(0)=\frac{1}{\beta_0}, \label{eq:E}
\end{equation}
так что $P=\cE(\delta)$, и $\beta(\delta)=v/c=pc/E=(1+\delta)/\cE(\delta)$, $\gamma=E/(mc^2)=\cE\,p_0/(mc)$.

\paragraph{Вывод гамильтониана по шагам.}
\emph{Шаг 0.} С независимой переменной $s$ (прямой участок) канонические пары $(x,P_x)$, $(y,P_y)$, $(t,-E)$,
а гамильтониан есть минус канонический продольный импульс:
\begin{equation}
K=-\sqrt{\frac{E^2}{c^2}-m^2c^2-(P_x-ZeA_x)^2-(P_y-ZeA_y)^2}\;-\;ZeA_s .
\end{equation}
\emph{Шаг 1: переход к $(q,P)$ и нормировка.} Производящая функция второго рода
$F_2(t,P;s)=-p_0c\,P\,\big(t-t_0(s)\big)$ даёт старый импульс $-E=\partial F_2/\partial t=-p_0cP$, новую координату
$\partial F_2/\partial P=-p_0c\,\Delta t=p_0\,q$ и добавку к гамильтониану $\partial F_2/\partial s=p_0cP/v_0=p_0P/\beta_0$.
Деля всё на $p_0$ (это сохраняет каноничность при одновременной нормировке импульсов), получаем
\begin{equation}
H(q,P)=\frac{P}{\beta_0}-\sqrt{P^2-\frac{1}{\beta_0^2\gamma_0^2}-(p_x-a_x)^2-(p_y-a_y)^2}-a_s,\quad
a_i=\frac{Ze}{p_0}A_i . \label{eq:HqP}
\end{equation}
\emph{Шаг 2: точечное преобразование $P\mapsto\delta$.} Функция $F_2(q,\delta)=q\,\cE(\delta)$ даёт
\begin{equation}
P=\frac{\partial F_2}{\partial q}=\cE(\delta),\qquad
l=\frac{\partial F_2}{\partial\delta}=q\,\frac{d\cE}{d\delta}=q\,\frac{1+\delta}{\cE}=q\,\beta(\delta),
\quad l=-\beta c\,\Delta t=-v\,\Delta t, \label{eq:point}
\end{equation}
и (так как $F_2$ не зависит от $s$) значение гамильтониана не меняется, достаточно подставить
$P=\cE(\delta)$, $q=l/\beta(\delta)$. Поскольку $P^2-1/(\beta_0^2\gamma_0^2)=(1+\delta)^2$,
\begin{equation}
\boxed{\ H=\frac{\cE(\delta)}{\beta_0}-p_s-a_s\Big(x,y,\ \frac{l}{\beta(\delta)};s\Big)\ },\qquad
p_s=\sqrt{(1+\delta)^2-(p_x-a_x)^2-(p_y-a_y)^2} \label{eq:Hld}
\end{equation}
Константу $\cE(0)/\beta_0=1/\beta_0^2$ можно вычесть. Проверка на дрейфе ($a=0$):
$dl/ds=\partial H/\partial\delta=\beta/\beta_0-(1+\delta)/p_s$, т.е. $dl/ds=v/v_0-1$ на оси,
как и должно быть для $l=-v\Delta t$ при $d\Delta t/ds=1/v-1/v_0$.

\section{Поле TM$_{010}$ и гамильтониан толстого резонатора}
В цилиндрическом резонаторе (мода TM$_{010}$, $k=\omega/c$) поле
\begin{equation}
E_z=E_s\,J_0(k\rho)\sin(\omega t+\varphi_0),\qquad B_\phi=\frac{E_s}{c}J_1(k\rho)\cos(\omega t+\varphi_0).
\end{equation}
Выбираем калибровку $\Phi=0$. Выбор калибровки свободен (физика калибровочно инвариантна: поля и траектории не
зависят от неё, меняется лишь вид гамильтониана на полную производную), а $\Phi=0$ здесь удобна, потому что поле
в резонаторе чисто вихревое ($\nabla\cdot\mathbf E=0$, свободных зарядов нет): $\mathbf E=-\partial\mathbf A/\partial t$
описывается одним векторным потенциалом. Тогда
\begin{equation}
A_z=\frac{E_s}{\omega}J_0(k\rho)\cos(\omega t+\varphi_0)
\end{equation}
(проверка: $-\partial A_z/\partial t=E_z$, $-\partial A_z/\partial\rho=B_\phi$). Подставляя
$\omega t=\omega\Delta t+\omega s/v_0=-kq+ks/\beta_0$ в \eqref{eq:Hld},
\begin{equation}
H_\rf=-a_s=-\frac{ZeE_s}{p_0\,\omega}\,J_0(k\rho)\cos\!\Big(\frac{ks}{\beta_0}-kq+\varphi_0\Big),\qquad q=\frac{l}{\beta(\delta)} .
\end{equation}
Зависимость от $s$ через $ks/\beta_0$ описывает пролёт частицы через осциллирующее поле. Усреднение по длине
резонатора $L=\pi/k$ (половина длины волны) даёт эффективное напряжение $V=E_sLT_{\rm tr}$ с фактором пролётного времени
\begin{equation}
T_{\rm tr}=\frac{\sin(\omega L/2v_0)}{\omega L/2v_0}
\end{equation}
($T_{\rm tr}=2/\pi$ для $L=\lambda/2$ и $\beta_0=1$), и набег энергии $ZeV\sin(\varphi_c-kq)$, где $\varphi_c$~--- фаза
поля в момент прохода центра резонатора опорной частицей. Вблизи оси $J_0(k\rho)\approx 1-k^2\rho^2/4$;
поперечный член (фокусировка ВЧ) порядка $k^2\rho^2$ далее опускается.

\paragraph{Модель однородного зазора.}
Обозначим $\varphi_s=\pi-\varphi_c$, так что $\sin(\varphi_c-kq)=\sin(\varphi_s+kq)$, и положим, что набег энергии
распределён по длине $L$ равномерно: $dE/ds=ZeV\sin(\varphi_s+kq)/L$ ($V$ уже включает $T_{\rm tr}$). Тогда
\begin{equation}
\frac{dP}{ds}=-\frac{\partial H_\rf}{\partial q}=\frac{ZeV}{p_0cL}\sin(\varphi_s+kq),\qquad
H_\rf=\frac{ZeV}{p_0c\,kL}\cos(\varphi_s+kq), \label{eq:Hrf}
\end{equation}
где $k=h\omega_0/c=2\pi h\beta_0/C_0$, $C_0$~--- длина опорной орбиты, $h$~--- кратность. Знак: $kq=-h\omega_0\Delta t$, т.е. частица, пришедшая
раньше ($\Delta t<0$), имеет $q>0$ и большую фазу $\varphi=\varphi_s+kq$. В переменных \eqref{eq:ld}
\begin{equation}
\boxed{\ H=\frac{\cE(\delta)}{\beta_0}-p_s+\frac{ZeV}{p_0c\,kL}\cos\!\Big(\varphi_s+\frac{k\,l}{\beta(\delta)}\Big)\ }
\label{eq:Hfull}
\end{equation}
(плюс члены поворота и мультиполей). Фаза ВЧ $\varphi-\varphi_s=k\,l/\beta(\delta)$; из-за $\delta$-зависимости $\beta$ член $H_\rf$ зависит и от $\delta$,
что даёт поправку к $dl/ds$: она в точности учитывает, что $l=-v\Delta t$ меняется при изменении скорости
при фиксированном $\Delta t$.

\section{Тонкий предел}
Отображение толстого резонатора за длину $L$ есть $\mathcal M=e^{-L:H:}$ с $H=H_d+H_\rf$, где $H_d=\cE/\beta_0-p_s$~---
дрейф, а $:f:g=[f,g]$~--- скобка Пуассона. Положим
$G=L\,H_\rf=\int_0^LH_\rf\,ds=\dfrac{ZeV}{p_0ck}\cos(\varphi_s+kq)$ и устремим $L\to0$ при фиксированном
напряжении $V$. Тогда показатель $-L:H:=-L:H_d:-:G:$: первое слагаемое пропорционально длине и исчезает
(за нулевую длину частица не дрейфует), а второе остаётся конечным, потому что при фиксированном $V$
амплитуда $H_\rf\propto1/L$ растёт. Значит, $\mathcal M\to e^{-:G:}$~--- чистый толчок. В паре $(q,P)$ функция $G$
зависит только от $q$, поэтому этот поток вычисляется точно:
\begin{equation}
q'=q,\qquad P'=P+\frac{ZeV}{p_0c}\sin(\varphi_s+kq),
\end{equation}
т.е. $\Delta t$ сохраняется, а энергия получает толчок $E'=E+ZeV\sin\varphi$. Переход к $(l,\delta)$ точный:
\begin{equation}
\cE'=\cE(\delta)+\frac{ZeV}{p_0c}\sin\varphi,\qquad \delta'=\sqrt{\cE'^2-\frac{1}{\beta_0^2\gamma_0^2}}-1,\qquad
l'=-v(\delta')\,\Delta t = l\,\frac{\beta(\delta')}{\beta(\delta)} . \label{eq:thin}
\end{equation}
Поперечные переменные не меняются.

\paragraph{Точность схемы <<дрейф--толчок--дрейф>>.}
Насколько толстый резонатор отличается от <<дрейф $L/2$ -- тонкий толчок $G$ -- дрейф $L/2$>>, показывает
формула Бейкера--Кэмпбелла--Хаусдорфа (БКХ). Для некоммутирующих операторов $X,Y$
\begin{equation}
e^{X}e^{Y}=\exp\!\Big(X+Y+\tfrac12[X,Y]+\tfrac1{12}\big([X,[X,Y]]+[Y,[Y,X]]\big)+\dots\Big),
\end{equation}
где многоточие~--- коммутаторы четвёртой и более высокой степени. Применяя её дважды к симметричному
произведению, получаем
\begin{equation}
e^{X/2}e^{Y}e^{X/2}=\exp\!\Big(X+Y-\tfrac1{24}[X,[X,Y]]-\tfrac1{12}[Y,[X,Y]]+\dots\Big):
\end{equation}
член $\tfrac12[X,Y]$ сокращается (вклады $\pm\tfrac14[X,Y]$ от левого и правого полудрейфов), и первые отличия
от $e^{X+Y}$~--- тройные коммутаторы (формулы проверены численно на случайных матрицах).
Положим $X=-L:H_d:$, $Y=-L:H_\rf:$ (для однородного зазора $H_\rf$ не зависит от $s$). Оба оператора
пропорциональны $L$, поэтому
\begin{equation}
e^{-\frac L2:H_d:}\,e^{-L:H_\rf:}\,e^{-\frac L2:H_d:}=e^{-L:(H_d+H_\rf):}\big(1+O(L^3)\big),
\end{equation}
тогда как несимметричная схема (<<дрейф $L$ -- толчок>>) имела бы ошибку $\tfrac12[X,Y]=O(L^2)$. Таким образом,
ошибка симметричной схемы на один резонатор~--- $O(L^3)$, и тонкий предел есть корректный предел той же
гамильтоновой системы. (Для реального поля $H_\rf$ зависит от $s$ через $ks/\beta_0$; усреднение по $s$ этот
порядок не меняет, а входит в $T_{\rm tr}$.)

\section{Симплектичность}
Отображение $e^{-:G:}$~--- поток гамильтоновой системы, следовательно симплектично в $(q,P)$, а точечное
преобразование \eqref{eq:point} каноническое; значит, толчок \eqref{eq:thin} симплектичен и в $(l,\delta)$:
$J^{\mathsf T}SJ=S$ для якобиана $J$ при любой амплитуде. Это свойство в нормированных переменных нарушается
только при смене опорного импульса (раздел~\ref{sec:acc}).

\section{Накопитель и ускорение}
\label{sec:acc}
\paragraph{Автономность.} Если гамильтониан периодичен по $s$ с периодом $C_0$ (одинаков на каждом обороте), то
система \emph{автономна} в смысле однооборотного отображения: одно и то же симплектическое отображение
$\mathcal M$ применяется на каждом обороте. Тогда определены замкнутая орбита (неподвижная точка $\mathcal M$),
линейные частоты (бетатронные и синхротронная $Q_s$), нормальная форма, инварианты и область устойчивости;
долговременная динамика~--- это итерации одного $\mathcal M$.

\paragraph{Накопитель.} Без ускорения ($\sin\varphi_s=0$, или $ZeV\sin\varphi_s=U_0$ ровно компенсирует потери на
излучение) опорная частица не набирает энергии, $p_0$ постоянно. Так как $kq$ отсчитывается от синхронной
частицы, а $h\omega_0T_0=2\pi h$, гамильтониан не зависит от номера оборота: система автономна в указанном
смысле, однооборотное отображение одно и то же и симплектично. Устойчивость: $\eta_1\cos\varphi_s>0$
($\varphi_s=0$ выше, $\varphi_s=\pi$ ниже критической энергии в этой записи через $\sin$).

\paragraph{Ускорение.} При $\sin\varphi_s\neq0$ в \eqref{eq:Hfull} есть линейный член
$-\dfrac{ZeV\sin\varphi_s}{p_0cL}\,q$: опорная орбита $l=\delta=0$ перестаёт быть решением. Естественная запись~---
переменный опорный импульс $p_0(s)$ с
\begin{equation}
\frac{dp_0}{ds}=\frac{ZeV\sin\varphi_s}{v_0L},
\end{equation}
т.е. синхронная частица остаётся опорной. Каноническое преобразование к переменным относительно $p_0(s)$~---
производящая функция $F_2=p_0(s)\,\big[x\,p_x+y\,p_y+l\,(1+\delta)\big]$ (старые импульсы размерные, новые
нормированы на $p_0(s)$); после деления на $p_0$ к $H$ добавляется
\begin{equation}
\Delta H=\frac{p_0'}{p_0}\Big[x\,p_x+y\,p_y+l\,(1+\delta)\Big]. \label{eq:dH}
\end{equation}
Проверка: $-\partial\Delta H/\partial l=-(p_0'/p_0)(1+\delta)$, а линейный член ВЧ даёт
$-\partial H_\rf/\partial l\big|_0=\dfrac{ZeV\sin\varphi_s}{\beta_0p_0cL}=\dfrac{1}{\beta_0p_0c}\dfrac{dE_0}{ds}=\dfrac{p_0'}{p_0}$
(так как $dE_0=v_0\,dp_0$), т.е. при $\delta=0$ они сокращаются и $l=\delta=0$ снова решение. Кроме того,
$\beta_0(s)$, $\gamma_0(s)$ в $\cE(\delta)/\beta_0$ и $k(s)=2\pi h\beta_0(s)/C_0$ становятся функциями $s$.

\paragraph{Неавтономность и потеря симплектичности.} При ускорении гамильтониан явно и непериодически зависит от
$s$ (через $p_0(s)$): система \emph{неавтономна}, и отображение за оборот $\mathcal M_n$ своё на каждом обороте.
Поток по-прежнему гамильтонов, поэтому в размерных (ненормированных) переменных каждое $\mathcal M_n$
симплектично. Но в нормированных переменных после резонатора делается скачок опорной энергии~--- перенормировка
(\cite{xsuite}, \S2.14, \S2.21); в наших переменных ($\Delta t$ и положение при этом не меняются)
\begin{equation}
l_{\rm new}=l_{\rm old},\qquad
1+\delta_{\rm new}=(1+\delta_{\rm old})\frac{p_{0,\rm old}}{p_{0,\rm new}},\qquad
p_{x,y,\rm new}=p_{x,y,\rm old}\frac{p_{0,\rm old}}{p_{0,\rm new}} .
\end{equation}
Определитель этой замены в каждой плоскости равен $p_{0,\rm old}/p_{0,\rm new}<1$ (в сумме
$(p_{0,\rm old}/p_{0,\rm new})^3$): она не симплектична, фазовый объём в нормированных переменных сжимается~---
это и есть адиабатическое затухание ($p_x\propto1/p_0$: сохраняется размерный $P_x$, а не $p_x$).

\paragraph{Адиабатический подход.} Если $p_0$ меняется медленно по сравнению с периодом синхротронных колебаний
($\dot\gamma_0/\gamma_0\ll Q_s\omega_0$), отображение <<замораживают>> при каждой энергии: при данном $p_0$ строят
автономное $\mathcal M(p_0)$, находят замкнутую орбиту, частоты и бакет, а затем медленно меняют $p_0$, считая,
что частица следует за этими мгновенными решениями (адиабатические инварианты сохраняются, амплитуды затухают
как степень $p_0$). Подход нарушается там, где параметры меняются быстро по сравнению с собственной частотой:
вблизи критической энергии $\gamma_t$ ($Q_s\to0$, раздел~\ref{sec:eom}) и при быстром прохождении резонансов.

\section{Частота синхротронных колебаний}
Линеаризуя однооборотное отображение (тонкий резонатор + кольцо с $\Delta l=-\eta_1C_0\,\delta$ за оборот, см.
\eqref{eq:map1}),
\begin{equation}
\Delta\delta=\frac{ZeV}{\beta_0^2E_0}\,\cos\varphi_s\,\frac{2\pi h}{C_0}\,l,\qquad
\cos 2\pi Q_s=1-\frac{2\pi h\,\eta_1\,ZeV\cos\varphi_s}{2\beta_0^2E_0},
\end{equation}
и при малом $Q_s$
\begin{equation}
Q_s=\sqrt{\frac{h\,\eta_1\,ZeV\cos\varphi_s}{2\pi\beta_0^2E_0}},\qquad \eta_1=\alpha_c-\frac{1}{\gamma_0^2}.
\end{equation}
Совпадение $Q_s$ из следа продольного блока однооборотной матрицы с этой формулой~--- проверка
знаков и нормировки \eqref{eq:Hrf} (отличие $O(Q_s^2)$).

\section{Продольные уравнения движения и однооборотное отображение}
\label{sec:eom}
\paragraph{Гамильтониан на оси.} Рассмотрим частицу на оси ($x=p_x=y=p_y=0$, $p_s=1+\delta$) в однородном зазоре
длины $L$. Из \eqref{eq:Hfull}
\begin{equation}
H(l,\delta)=\frac{\cE(\delta)}{\beta_0}-(1+\delta)+\frac{ZeV}{p_0c\,kL}\cos\varphi,\qquad
\varphi=\varphi_s+\frac{k\,l}{\beta(\delta)},\qquad k=\frac{2\pi h\beta_0}{C_0}. \label{eq:Hax}
\end{equation}
(Если в фазе заменить $l/\beta(\delta)$ на $l/\beta_0$, теряется член $\propto kl\,\delta/\gamma_0^2$; ниже всё точно.)

\paragraph{Уравнение для $l$.} $dl/ds=\partial H/\partial\delta$ по слагаемым:
\begin{itemize}
\item $\dfrac{\partial}{\partial\delta}\dfrac{\cE}{\beta_0}=\dfrac{1}{\beta_0}\dfrac{1+\delta}{\cE}=\dfrac{\beta}{\beta_0}=\dfrac{v}{v_0}$;
\quad $\dfrac{\partial}{\partial\delta}\big[-(1+\delta)\big]=-1$;
\item $\dfrac{\partial\varphi}{\partial\delta}=-\dfrac{k\,l}{\beta^2}\dfrac{d\beta}{d\delta}$, причём
$\dfrac{d\beta}{d\delta}=\dfrac{d}{d\delta}\dfrac{1+\delta}{\cE}=\dfrac{1}{\cE^3\beta_0^2\gamma_0^2}=\dfrac{\beta}{(1+\delta)\gamma^2}$,
поэтому $\dfrac{\partial}{\partial\delta}\cos\varphi=\dfrac{k\,l\,\sin\varphi}{\beta\,(1+\delta)\gamma^2}$.
\end{itemize}
\paragraph{Уравнение для $\delta$.} $d\delta/ds=-\partial H/\partial l$; от $l$ зависит только фаза,
$\partial\varphi/\partial l=k/\beta$, поэтому $-\partial_l\big[\tfrac{ZeV}{p_0ckL}\cos\varphi\big]=\tfrac{ZeV}{p_0cL}\tfrac{\sin\varphi}{\beta}$.

\paragraph{Система в $(l,\delta)$.} Итого
\begin{equation}
\boxed{\ \left\{\begin{aligned}
\frac{dl}{ds}&=\frac{\beta(\delta)}{\beta_0}-1+\frac{ZeV}{p_0c\,L}\,\frac{l\,\sin\varphi}{\beta\,(1+\delta)\,\gamma^2},\\[2pt]
\frac{d\delta}{ds}&=\frac{ZeV}{p_0c\,L}\,\frac{\sin\varphi}{\beta(\delta)},
\end{aligned}\right.\qquad \varphi=\varphi_s+\frac{k\,l}{\beta(\delta)}\ } \label{eq:eomld}
\end{equation}
Второе уравнение есть $dp/ds=(1/v)\,dE/ds$ с $dE/ds=ZeV\sin\varphi/L$; последний член первого есть
$q\,d\beta/ds$, т.е. изменение $l=\beta q$ за счёт изменения скорости при постоянном $\Delta t$.

\paragraph{Переменные $(\varphi,\Delta E)$.} Фаза прибытия и отклонение энергии
\begin{equation}
\varphi=\varphi_s+kq=\varphi_s+\frac{2\pi h}{C_0}\frac{\beta_0}{\beta(\delta)}\,l,\qquad
\Delta E=E-E_0=p_0c\Big(\cE(\delta)-\frac1{\beta_0}\Big); \label{eq:trans}
\end{equation}
обратно
\begin{equation}
\delta=\sqrt{\Big(\frac1{\beta_0}+\frac{\Delta E}{p_0c}\Big)^2-\frac1{\beta_0^2\gamma_0^2}}-1,\qquad
l=\frac{C_0}{2\pi h}\frac{\beta(\delta)}{\beta_0}\,(\varphi-\varphi_s).
\end{equation}
Пара $(\varphi,\Delta E)$~--- это $(q,P)$ с постоянными множителями и сдвигами ($\varphi-\varphi_s=kq$,
$\Delta E=p_0c\,P-E_0$), поэтому она канонична с точностью до постоянного множителя $kp_0c$ при гамильтониане.
Из \eqref{eq:HqP} ($dq/ds=1/\beta_0-1/\beta$, $dP/ds=ZeV\sin\varphi/(p_0cL)$):
\begin{equation}
\left\{\begin{aligned}
\frac{d\varphi}{ds}&=\frac{2\pi h}{C_0}\Big(1-\frac{\beta_0}{\beta(\delta)}\Big),\\[2pt]
\frac{d\Delta E}{ds}&=\frac{ZeV}{L}\,\sin\varphi .
\end{aligned}\right. \label{eq:eomphi}
\end{equation}
В этих переменных энергетический толчок линеен и точен, а вся нелинейность $\delta(\Delta E)$ уходит в фазу;
в $(l,\delta)$ наоборот, толчок нелинеен через $\cE(\delta)$.

\paragraph{Время оборота.} Пусть кольцо (без ВЧ) переводит частицу с отклонением $\delta$ по замкнутой орбите
длины $C(\delta)$ со скоростью $v(\delta)$ за время $T(\delta)=C(\delta)/v(\delta)$; $T_0=C_0/v_0$, $h\omega_0T_0=2\pi h$.
Два вклада:
\begin{itemize}
\item \emph{длина пути}: $x=D_x\delta+\tfrac12D_x^{(2)}\delta^2+\dots$, $C(\delta)=C_0+\oint x/\rho\,ds+\dots$, т.е.
$C/C_0=1+\alpha_c\delta+\alpha_1\delta^2+O(\delta^3)$, $\alpha_c=\frac{1}{C_0}\oint\frac{D_x}{\rho}ds$
(в $\alpha_1$ входят $D_x^{(2)}$ и $D_x'^2$);
\item \emph{скорость}: $\beta(\delta)/\beta_0=(1+\delta)\big/\sqrt{1+\beta_0^2(2\delta+\delta^2)}$,
$\ \beta_0/\beta=1-\delta/\gamma_0^2+(1+\beta_0^2/2)\,\delta^2/\gamma_0^2+\dots$
\end{itemize}
Перемножая,
\begin{equation}
\frac{T(\delta)}{T_0}-1=\frac{C(\delta)}{C_0}\frac{\beta_0}{\beta(\delta)}-1=\eta_1\delta+\eta_2\delta^2+O(\delta^3),
\end{equation}
\begin{equation}
\boxed{\ \eta_1=\alpha_c-\frac{1}{\gamma_0^2},\qquad
\eta_2=\alpha_1-\frac{\alpha_c}{\gamma_0^2}+\frac{3\beta_0^2}{2\gamma_0^2}+\frac{1}{\gamma_0^4}
=\alpha_1-\frac{\eta_1}{\gamma_0^2}+\frac{3\beta_0^2}{2\gamma_0^2}\ } \label{eq:eta2}
\end{equation}
(проверено символьно). При $\gamma_0\to\infty$: $\eta_2\to\alpha_1$. За оборот без ВЧ $\Delta t$ растёт на $T-T_0$,
поэтому $\Delta q=-c(T-T_0)$ и $\Delta l=-v(T-T_0)=C_0\,\beta/\beta_0-C(\delta)$ (точно).

\paragraph{Однооборотное отображение.} Резонатор (тонкий, точный толчок \eqref{eq:thin}) и затем кольцо. Обозначим
\[
S_n=\sin\varphi_n-\sin\varphi_s,\qquad A=\frac{ZeV}{\beta_0^2E_0},\qquad \frac{ZeV}{p_0c}=\beta_0A,\]
\[\varepsilon=\frac{\Delta E}{\beta_0^2E_0},\qquad \beta_n=\beta(\delta_n),\quad T_n=T(\delta_n).
\]
Член $-\sin\varphi_s$ описывает либо потери на излучение $U_0=ZeV\sin\varphi_s$ (накопитель), либо, при ускорении,
переход к опорному импульсу $p_0(s)$ (раздел~\ref{sec:acc}). Ниже формулы в двух наборах переменных записаны
в одинаковом порядке: точно, первый порядок, второй порядок по $(\delta,\,l)$ или $(\varepsilon,\,\varphi-\varphi_s)$.

\medskip\noindent\textbf{Точно.}
\begin{align}
(l,\delta):\quad &\cE_{n+1}=\cE(\delta_n)+\beta_0A\,S_n,\qquad \delta_{n+1}=\sqrt{\cE_{n+1}^2-\tfrac{1}{\beta_0^2\gamma_0^2}}-1, \notag\\
&l_{n+1}=\frac{\beta_{n+1}}{\beta_n}\,l_n-C_0\,\frac{\beta_{n+1}}{\beta_0}\Big(\frac{T_{n+1}}{T_0}-1\Big),
\label{eq:mapex_l}\\[4pt]
(\varphi,\Delta E):\quad &\Delta E_{n+1}=\Delta E_n+ZeV\,S_n, \notag\\
&\varphi_{n+1}=\varphi_n-2\pi h\Big(\frac{T_{n+1}}{T_0}-1\Big),
\label{eq:mapex_phi}
\end{align}
где в $(l,\delta)$ фаза $\varphi_n=\varphi_s+\frac{2\pi h}{C_0}\frac{\beta_0}{\beta_n}l_n$, а в $(\varphi,\Delta E)$ $\delta_{n+1}=\delta(\Delta E_{n+1})$
по \eqref{eq:trans}. Второе слагаемое в $l_{n+1}$ равно $C_0\beta_{n+1}/\beta_0-C(\delta_{n+1})$; первое~--- пересчёт
$l=\beta q$ при толчке ($\Delta t$ не меняется). Деля $l_{n+1}$ на $\beta_{n+1}/\beta_0$ и умножая на $2\pi h/C_0$,
получаем ровно \eqref{eq:mapex_phi}: это одно и то же отображение.

\medskip\noindent\textbf{Первый порядок.}
\begin{align}
(l,\delta):\quad &\delta_{n+1}=\delta_n+A\,S_n,\qquad \varphi_n\simeq\varphi_s+\tfrac{2\pi h}{C_0}l_n, \notag\\
&l_{n+1}=l_n-C_0\,\eta_1\,\delta_{n+1}, \label{eq:map1}\\[4pt]
(\varphi,\Delta E):\quad &\varepsilon_{n+1}=\varepsilon_n+A\,S_n, \notag\\
&\varphi_{n+1}=\varphi_n-2\pi h\,\eta_1\,\varepsilon_{n+1}. \label{eq:map1phi}
\end{align}

\medskip\noindent\textbf{Второй порядок.}
\begin{align}
(l,\delta):\quad &\delta_{n+1}=\delta_n+A\,S_n\Big(1-\frac{\delta_n}{\gamma_0^2}\Big)-\frac{A^2S_n^2}{2\gamma_0^2},\qquad
\varphi_n\simeq\varphi_s+\tfrac{2\pi h}{C_0}\,l_n\Big(1-\frac{\delta_n}{\gamma_0^2}\Big), \notag\\
&l_{n+1}=\Big(1+\frac{\delta_{n+1}-\delta_n}{\gamma_0^2}\Big)l_n-C_0\Big[\eta_1\delta_{n+1}+\Big(\eta_2+\frac{\eta_1}{\gamma_0^2}\Big)\delta_{n+1}^2\Big],
\label{eq:map2}\\[4pt]
(\varphi,\Delta E):\quad &\varepsilon_{n+1}=\varepsilon_n+A\,S_n, \notag\\
&\varphi_{n+1}=\varphi_n-2\pi h\Big[\eta_1\varepsilon_{n+1}+\Big(\eta_2-\frac{\eta_1}{2\gamma_0^2}\Big)\varepsilon_{n+1}^2\Big]. \label{eq:map2phi}
\end{align}
Здесь использованы разложения (проверены символьно)
$\delta=\varepsilon-\varepsilon^2/(2\gamma_0^2)+\dots$,
$\delta'=\delta+\tfrac{\epsilon}{\beta_0}\big(1-\tfrac{\delta}{\gamma_0^2}\big)-\tfrac{\epsilon^2}{2\beta_0^2\gamma_0^2}$ для толчка $\cE\to\cE+\epsilon$,
$\beta_{n+1}/\beta_n\simeq1+(\delta_{n+1}-\delta_n)/\gamma_0^2$ и
$(\beta/\beta_0)(T/T_0-1)=\eta_1\delta+(\eta_2+\eta_1/\gamma_0^2)\delta^2+\dots$
В смешанной записи $(\varphi,\delta)$ (фаза из \eqref{eq:mapex_phi}, импульс из \eqref{eq:mapex_l}):
$\varphi_{n+1}=\varphi_n-2\pi h(\eta_1\delta_{n+1}+\eta_2\delta_{n+1}^2)$. Коэффициент при квадрате зависит от
выбора переменной ($\eta_2$ для $\delta$ в фазе, $\eta_2+\eta_1/\gamma_0^2$ для $l$, $\eta_2-\eta_1/(2\gamma_0^2)$ для
$\varepsilon$), но нули $T(\delta)=T_0$ у всех трёх одни и те же.

\paragraph{Следствия.} Для медленных колебаний ($Q_s\ll1$) номер оборота $n$ можно считать непрерывным. В смешанной
записи
\begin{equation}
\frac{d\varphi}{dn}=-2\pi h(\eta_1\delta+\eta_2\delta^2),\qquad \frac{d\delta}{dn}=A(\sin\varphi-\sin\varphi_s),
\end{equation}
с гамильтонианом (по $n$)
\begin{equation}
K(\varphi,\delta)=-2\pi h\Big(\frac{\eta_1\delta^2}{2}+\frac{\eta_2\delta^3}{3}\Big)+A\big(\cos\varphi+\varphi\sin\varphi_s\big). \label{eq:K}
\end{equation}
Тот же гамильтониан в канонических $(l,\delta)$ (с $\varphi-\varphi_s\simeq2\pi h\,l/C_0$):
\begin{equation}
K_l(l,\delta)=-C_0\Big(\frac{\eta_1\delta^2}{2}+\frac{\tilde\eta_2\delta^3}{3}\Big)
+\frac{AC_0}{2\pi h}\Big[\cos\Big(\varphi_s+\frac{2\pi h\,l}{C_0}\Big)+\frac{2\pi h\,l}{C_0}\sin\varphi_s\Big]
\end{equation}
с $\tilde\eta_2=\eta_2+\eta_1/\gamma_0^2$,
$dl/dn=\partial K_l/\partial\delta$, $d\delta/dn=-\partial K_l/\partial l$.
\begin{itemize}
\item \emph{Только $\eta_1$}: маятник. Неподвижные точки $\delta=0$, $\sin\varphi=\sin\varphi_s$; точка $\varphi_s$
устойчива при $\eta_1\cos\varphi_s>0$, линеаризация даёт $Q_s^2=hA\eta_1\cos\varphi_s/(2\pi)$, как в разделе~6.
Сепаратриса~--- траектория через неустойчивую точку~--- ограничивает область устойчивых колебаний, которую
принято называть \emph{бакетом} (англ.\ \emph{bucket}, <<ковш>>); ниже используется этот термин. Сепаратриса
симметрична по $\delta\to-\delta$; для $\sin\varphi_s=0$ полувысота бакета $\delta_{\max}=2\sqrt{A/(\pi h|\eta_1|)}$.
\item \emph{С $\eta_2$}: $d\varphi/dn=0$ также при $\delta_*=-\eta_1/\eta_2$. Появляется вторая цепочка неподвижных точек
$(\varphi,\delta_*)$ с $\sin\varphi=\sin\varphi_s$, на которой эффективный коэффициент проскальзывания
$\partial(\eta_1\delta+\eta_2\delta^2)/\partial\delta=-\eta_1$ имеет противоположный знак: устойчивая точка при
$\delta=0$ становится неустойчивой при $\delta_*$ и наоборот (``альфа-бакеты''). Кубический член в $K$ делает
сепаратрису несимметричной по $\delta$: бакет сжат со стороны $\delta_*$ и растянут с другой.
Пока $|\delta_*|\gg\delta_{\max}$, поправка мала, $\propto\eta_2\delta_{\max}/\eta_1$.
\item \emph{Вблизи $\gamma_t$} ($\alpha_c=1/\gamma_t^2$, $\eta_1\to0$): $Q_s\propto\sqrt{|\eta_1|}\to0$, адиабатичность нарушается,
$\delta_{\max}\propto|\eta_1|^{-1/2}$ растёт, и $\delta_*=-\eta_1/\eta_2\to0$ попадает внутрь бакета~---
нелинейный член $\eta_2\delta^2$ определяет динамику (асимметрия, потери частиц при прохождении
критической энергии). При $\eta_1=0$ остаётся $d\varphi/dn=-2\pi h\eta_2\delta^2$, и линейная фокусировка отсутствует.
\end{itemize}

\section{Итог}
\begin{itemize}
\item Гамильтониан \eqref{eq:Hfull} выведен прямо в $(l,\delta)$; ВЧ-резонатор описывается так же, как гамильтонианы
обычных элементов; толстый резонатор симметрично разбивается на дрейф--толчок--дрейф с ошибкой $O(L^3)$.
\item Накопитель: $p_0=\mathrm{const}$, система автономна, однооборотное отображение одно и симплектично.
\item Ускорение: опорный импульс $p_0(s)$, система неавтономна; перенормировка после резонатора несимплектична в
нормированных переменных~--- адиабатическое затухание.
\item Однооборотное отображение \eqref{eq:mapex_l}--\eqref{eq:map2phi} дано точно и с разложением до $\eta_2$ в обоих
наборах переменных.
\end{itemize}

\appendix
\section{Переменные $(z,\delta_E)$ и пересчёт коэффициентов}
В лекциях \cite{wolski4,wolski7} используется продольная пара $(z,\delta_E)$: $z=s/\beta_0-ct=q$ и
\emph{энергетическое} отклонение $\delta_E=E/(p_0c)-1/\beta_0=P-1/\beta_0=\Delta E/(p_0c)$~--- то же, что $(q,P)$
со сдвигом импульса, т.е. каноническая пара; $\varepsilon=\delta_E/\beta_0$. Связь с нашими:
$1+\delta=\sqrt{(\delta_E+1/\beta_0)^2-1/(\beta_0^2\gamma_0^2)}$, $l=\beta(\delta)\,z$, $d\delta_E=\beta\,d\delta$.
Там длина пути и период разлагаются по $\delta_E$:
$C/C_0=1+\alpha_p\delta_E+\tfrac12\alpha_p^{(2)}\delta_E^2$, $T/T_0=1+\eta_p\delta_E+\tfrac12\eta_p^{(2)}\delta_E^2$,
а при $\nu_s\ll1$
\[
\frac{dz}{ds}\approx-\eta_p\,\delta_E,\qquad
\frac{d\delta_E}{ds}\approx\frac{1}{C_0}\frac{ZeV}{p_0c}\big[\sin(\varphi_s^{(z)}-kz)-\sin\varphi_s^{(z)}\big],
\qquad \varphi_s^{(z)}=\pi-\varphi_s,
\]
что совпадает с \eqref{eq:map1phi} (с $s=nC_0$, $\beta_0\to1$). Подставляя
$\delta=\delta_E/\beta_0-\delta_E^2/(2\beta_0^2\gamma_0^2)+O(\delta_E^3)$ в $\eta_1\delta+\eta_2\delta^2$:
\[
\eta_p=\frac{\eta_1}{\beta_0},\quad \tfrac12\eta_p^{(2)}=\frac{\eta_2}{\beta_0^2}-\frac{\eta_1}{2\beta_0^2\gamma_0^2},\qquad
\alpha_p=\frac{\alpha_c}{\beta_0},\quad \tfrac12\alpha_p^{(2)}=\frac{\alpha_1}{\beta_0^2}-\frac{\alpha_c}{2\beta_0^2\gamma_0^2}.
\]
\paragraph{Замечание об $\eta_p$.} Отсюда $\eta_p=\alpha_p-1/(\beta_0\gamma_0^2)$ и $\gamma_t^2=1/\alpha_c=1/(\beta_0\alpha_p)$.
В \cite{wolski7} (ур.~34--35) стоит $\eta_p=\alpha_p-1/\gamma_0^2$, $\gamma_t^2=1/\alpha_p$~--- это верно при $\beta_0\to1$
(ультрарелятивистский случай, который там и подразумевается). Множитель $\beta_0$ появляется потому, что $\alpha_p$
определено по энергии: $[(1/\beta)\,d\beta/d\delta_E]_0=1/(\beta_0\gamma_0^2)$. (В ур.~31 там, по-видимому, опечатка:
напечатано $\beta\gamma$ вместо $\gamma=(1/\beta_0+\delta_E)\beta_0\gamma_0$.) В импульсной форме
$\eta_1=\alpha_c-1/\gamma_0^2$ множителя нет, и $\gamma_t^2=1/\alpha_c$ точно.

\begin{thebibliography}{9}
\bibitem{wolski4} A.~Wolski, \emph{Linear and nonlinear dynamics, Lecture 4: RF cavities and acceleration},
Cockcroft Institute (2014), \url{https://www.cockcroft.ac.uk/wp-content/uploads/2014/12/wolski-4.pdf}.
\bibitem{wolski7} A.~Wolski, \emph{Lecture 7: Longitudinal dynamics}, Cockcroft Institute (2014).
\bibitem{wolskibook} A.~Wolski, \emph{Beam Dynamics in High Energy Particle Accelerators}, Imperial College Press (2014).
\bibitem{tecker} F.~Tecker, \emph{Longitudinal beam dynamics}, CAS -- CERN Accelerator School:
Introduction to Accelerator Physics (2020), arXiv:2011.02932.
\bibitem{xsuite} Xsuite physics manual, \S2.14 (Cavity), \S2.21 (Reference energy increase),
\url{https://xsuite.readthedocs.io}.
\bibitem{leduff} J.~Le~Duff, \emph{Longitudinal beam dynamics in circular accelerators}, CAS, CERN 94-01 (1994).
\end{thebibliography}
\end{document}
